% !TeX root = ../main.tex
% 系统概要设计

\chapter{系统概要设计}
    在软件架构设计领域，Kruchten \cite{Kruchten1995} 提出的 \textbf{4+1 视图模型} 已成为描述复杂软件系统架构的标准方法。该模型通过逻辑视图（Logical View）、开发视图（Development View）、过程视图（Process View）、物理视图（Physical View）以及场景视图（Scenarios / Use-Case View），从多个维度对系统架构进行全面描述，以满足系统开发、实现、测试与运维等不同角色的关注点。

    本文的系统概要设计采用 4+1 视图方法进行组织。鉴于前文需求分析阶段已通过用例设计实现了场景视图，本章将重点围绕逻辑视图、开发视图、过程视图和物理视图展开，系统化、结构化地呈现系统的模块划分、运行机制、实现组织及部署架构。

\section{逻辑视图}
    逻辑视图（Logical View）面向系统逻辑分析和设计，描述系统逻辑结构的视图，主要解决系统分析和设计的问题，它描述系统的业务上下文、系统的逻辑分解，以及分解出的逻辑元素间的关系。本节通过逻辑模型和技术模型来描述逻辑视图。

\subsection{逻辑模型}
    逻辑模型描述系统的逻辑功能模块分解，将系统分解为相应的逻辑功能元素，并描述各逻辑功能元素之间的关系。以下通过系统架构设计与功能模块分解来构造逻辑模型。

\subsubsection{系统架构设计}
    系统逻辑架构如图~\ref{fig:4.1}所示，整体分为四个主要层级：配置管理层、数据持久层、业务处理层和用户展示层。这些层级之间的交互接口也进行了详细设计，具体内容如下：
    \begin{itemize}
        \item \textbf{配置管理层}：该层负责集中管理系统的各类配置信息，包括 SmartCI 配置管理和系统服务配置管理。SmartCI 配置管理组件通过 XML 配置文件定义任务单元和流水线，确保系统集成到 SmartCI 平台。而系统服务配置管理则包含数据库配置、Flask 服务配置等。
        \item \textbf{数据持久层}：该层包含数据库模型和静态文件系统。数据库模型负责存储覆盖率数据及相关配置，通过 Flask 框架的 Flask-Migrate 执行数据库迁移，自动创建所需数据表。静态文件系统则利用 Flask 框架中的 static 目录存储覆盖率报告，用户可以通过构造 URL 直接访问这些报告文件。
        \item \textbf{业务处理层}：该层由覆盖率采集组件和覆盖率解析组件构成。覆盖率解析组件是系统的核心模块，基于 Flask 框架实现，使用 Flask-RESTful 创建接口，提供上传覆盖率文件、查询解析结果、配置及查询解析参数等功能。覆盖率采集组件则负责从冒烟测试中收集覆盖率数据，并调用接口将数据上传至覆盖率解析组件进行进一步处理。
        \item \textbf{用户展示层}：这一层采用 Vue.js 框架构建，包含页面组件和 Vuex 状态管理组件。Vuex 负责全局状态管理，并调用后端接口与后端交互，同时构造页面接口将其传递给页面组件。页面组件接收数据后，将覆盖率数据可视化呈现。
    \end{itemize}

    \begin{figure}[ht]
        \centering
        \includegraphics[width=0.9\textwidth]{logical/逻辑视图-系统架构图.pdf}
        \caption{逻辑视图-系统架构图}
        \label{fig:4.1}
    \end{figure}

\subsubsection{功能模块划分}
    功能模块图如图~\ref{fig:4.2}所示，可分为以下几个部分：

    \begin{figure}[ht]
        \centering
        \includegraphics[width=0.8\textwidth]{logical/逻辑视图-功能模块图.pdf}
        \caption{逻辑视图-功能模块图}
        \label{fig:4.2}
    \end{figure}

    \begin{itemize}
        \item \textbf{覆盖率采集模块}：从冒烟测试流水线中收集覆盖率数据，支持多种交叉编译工具链（LLVM 和 GCC）。该模块与冒烟测试流水线集成，确保覆盖率数据的实时采集。
        \item \textbf{细粒度分析模块}：处理采集到的覆盖率数据，支持按单板、领域、模块等多个维度进行分层统计，并计算增量覆盖率，同时提供多时间点的覆盖率对比与趋势分析功能。
        \item \textbf{覆盖率可视化模块}：通过前端界面向用户展示分层次、多维度、细粒度的覆盖率数据视图，包括柱状图、覆盖率报告等。
        \item \textbf{数据存储模块}：将解析得到的覆盖率数据持久化存储到数据库中，并将生成的覆盖率报告保存至静态文件系统，确保数据的长期可追溯和可复用。
    \end{itemize}

\subsection{技术模型}
    技术模型定义系统采用的关键技术部件和技术栈，通常包括整体框架技术，公共机制，基础设施，公共服务/组件，以及各逻辑功能元素的技术方案等。

    系统技术模型如图~\ref{fig:4.3}所示，描述了系统各个逻辑组件使用的重要技术。具体内容如下：

    \begin{itemize}
        \item 整个系统使用 Flask 框架实现。
        \item 用户展示层通过 Vue.js 框架实现，其中页面组件使用 Element Plus 组件库 和 Echarts 组件快速构建页面元素，Vuex 管理组件使用 Vue.js 框架中的 Vuex 工具实现全局状态管理。
        \item 业务处理层中，覆盖率解析组件主要使用 lcov 工具，将原始的覆盖率文件处理为标准格式，并执行数据筛选、细粒度划分与增量覆盖率提取操作,实现细粒度覆盖率分析。覆盖率采集组件集成到冒烟测试流水线中，通过LLVM/GCC工具链插桩采集覆盖率数据，并使用共享内存通信中间件进行跨OS通信实现覆盖率完整采集。
        \item 数据持久层中，数据库Model组件使用了MySQL数据库来存储覆盖率数据，静态文件系统组件依赖于 Flask 框架提供的静态文件服务，通过 Flask 的 static 目录存储覆盖率报告，用户可以通过形如 http://ip:port/static/ 的 URL 访问存储的静态资源。
    \end{itemize}
    
    \begin{figure}[ht]
        \centering
        \includegraphics[width=1.0\textwidth]{logical/逻辑视图-技术模型.pdf}
        \caption{逻辑视图-技术模型图}
        \label{fig:4.3}
    \end{figure}

\section{开发视图}
    开发视图（Development View）面向系统开发及软件管理，描述系统代码结构，构建结构的视图，主要解决系统技术实现和开发的问题，它依托逻辑视图，描述代码、构建结构。本节通过代码模型来描述开发视图。

    
\subsection{代码模型}
    代码模型图如图~\ref{fig:4.4}所示，定义代码结构以及代码元素与逻辑模型中逻辑元素的对应关系，建立逻辑元素到代码目录的映射关系，以实现软件源代码的显示管理。

    \begin{figure}[ht]
        \centering
        \includegraphics[width=1.0\textwidth]{develop/开发视图-代码模型.pdf}
        \caption{开发视图-代码模型图}
        \label{fig:4.4}
    \end{figure}

    核心代码存储在代码库 \verb|BspCoverage| 中，由 Flask 框架实现。

    其中，\verb|/src| 目录包含后端代码实现，其中每个子目录代表一个独立的应用。在代码模型中，覆盖率解析组件和数据库 Model 组件位于 \verb|/src/coverage| 目录下，二者归属于同一个应用。

    静态文件服务依赖于 Flask 框架的特性，静态资源存储在 \verb|/static| 目录中，用户可以通过形如 \verb|http://ip:port/static/file| 的 URL 访问存储在该目录下文件，其中 \verb|ip| 和 \verb|port| 代表 Flask 应用所在服务器的地址与端口，\verb|file| 是需要访问的静态文件名。

    \verb|/web| 目录包含了前端代码实现，采用 Vue.js 框架进行开发，并使用 Vuex 进行全局状态管理。具体地，\verb|/src/store| 目录用于存放状态管理代码，\verb|/src/components| 存放前端组件，\verb|/src/views| 存放页面实现。每个页面可能由多个组件构成，并通过 Vue.js 渲染展示给用户。

    通过\verb|/.config.py| 文件实现系统服务配置管理，通过 \verb|/.smartci| 目录下的 \verb|task.xml| 和 \verb|project.xml| 实现 SmartCI 配置管理。

    注意到，在代码模型中，覆盖率解析组件依赖于覆盖率采集组件，但后者的实现并未包含在当前代码库中。本系统的覆盖率采集组件基于 C 代码实现，主要实现了基于共享内存的跨 OS 通信中间件，负责覆盖率的完整采集与传输。同时，完整的覆盖率采集流程由项目组其他成员负责实现，集成到冒烟测试流水线中，因此当前代码库不包含覆盖率采集组件。

\section{过程视图}
    过程视图（Process View）又称运行视图，面向系统运行，描述系统启动过程、运行期交互的视图，主要解决系统运行期交互问题，描述各可执行交付件在运行期的交互关系。本节通过各个模块的运行顺序图来描述过程视图。

\subsection{系统总体流程}
    首先，在用户视角下，绘制系统总体的过程视图，如图~\ref{fig:4.5}所示。该时序图主要描述了从用户交互到数据存储和可视化呈现的系统运行时序。整个流程可概括为三个阶段：解析参数配置、冒烟测试与覆盖率分析，以及覆盖率检视。

    \begin{figure}[ht]
        \centering
        \includegraphics[width=0.97\textwidth]{process/过程视图-覆盖率看板总体流程.pdf}
        \caption{过程视图-系统总体流程时序图}
        \label{fig:4.5}
    \end{figure}

    在解析参数配置阶段，用户通过前端系统界面配置覆盖率解析所需的参数。前端系统将用户输入的参数通过接口提交给后端系统，后端系统对参数进行验证并存储至数据库中，同时接收数据库返回的执行结果以确认参数是否存储成功，并将操作结果反馈至前端界面，在用户界面显示配置成功或失败的提示信息。

    进入冒烟测试与覆盖率分析阶段，用户编写测试用例提交到测试环境中，触发冒烟测试流水线。测试环境执行交叉编译操作生成构建包和可执行文件，执行可执行文件生成覆盖率数据文件。将生成的文件返回用户主机，由用户上传至后端系统。其中构建包可选，是由于在 GCC 工具链下，无需上传构建包。 后端系统接收到上传文件后，首先从数据库获取解析参数，然后执行覆盖率数据解析，生成覆盖率数据与覆盖率报告。其中覆盖率数据存储到数据库中，覆盖率报告写入静态文件系统。解析完成后，后端系统通过接口返回响应信息，用户可以查看上传文件后的覆盖率解析结果。

    在覆盖率检视阶段，用户通过前端界面请求查看覆盖率分析结果。前端系统调用后端接口获取解析后的覆盖率数据，并通过静态文件系统访问已生成的覆盖率报告。后端系统从数据库检索覆盖率数据后返回前端，前端将数据与报告整合，构建交互式可视化界面，包括表格、柱状图以及告警信息等，并最终呈现给用户。通过这一阶段，用户能够直观地了解覆盖率分布和质量状况，实现覆盖率数据的可视化呈现。    

\subsection{覆盖率采集流程}
    覆盖率采集模块的过程视图如图~\ref{fig:4.6}所示，完整描述了用户、编译工具链及目标硬件平台各组件之间的交互过程。

    \begin{figure}[ht]
        \centering
        \includegraphics[width=0.9\textwidth]{process/过程视图-覆盖率采集流程.pdf}
        \caption{过程视图-覆盖率采集流程时序图}
        \label{fig:4.6}
    \end{figure}

    首先，用户在用户主机上准备交叉编译工具链并编写测试用例。用户主机将测试用例提交给交叉编译工具链进行编译，同时启用覆盖率选项，生成构建包及可执行文件。编译完成后，构建产物返回用户主机，形成下一步执行的输入。

    随后，用户主机将可执行文件分别分发至目标硬件平台的控制面（LinuxOS）和数据面（异构OS）。数据面在接收可执行文件后立即执行测试用例，并在执行过程中通过插桩技术采集覆盖率信息。采集到的统计数据通过共享内存通信机制实时传输至控制面，以保证在数据面缺乏独立文件系统且可能自动断电的环境下，实现覆盖率数据的完整采集。控制面在自身执行测试用例并采集覆盖率信息的同时，整合来自数据面的统计数据，生成统一的覆盖率数据文件。最终，控制面将汇总后的覆盖率数据文件返回至用户主机，为后续覆盖率解析和可视化处理提供可靠的基础数据。

\subsection{覆盖率解析流程}
    覆盖率解析模块的过程视图如图~\ref{fig:4.7}所示，展示了系统接收原始覆盖率数据、调用外部工具进行分析处理，并最终将结构化覆盖率数据存储到数据库的完整过程。
    
    \begin{figure}[ht]
        \centering
        \includegraphics[width=0.9\textwidth]{process/过程视图-覆盖率解析流程.pdf}
        \caption{过程视图-覆盖率解析流程时序图}
        \label{fig:4.7}
    \end{figure}

    LLVM 编译器与 GCC 编译器在生成和处理覆盖率数据时存在显著差异，这直接影响系统对覆盖率信息的解析方式。
    \begin{itemize}
        \item GCC 工具链在编译阶段会生成 .gcno 文件，包含程序控制流图及源代码行映射信息，运行时生成的 .gcda 文件仅记录各代码行的执行次数。利用 .gcno 和 .gcda 文件，解析工具可以直接重建完整覆盖率信息。而 .gcno 文件会自动包含在覆盖率采集模块输出的覆盖率数据文件中。
        \item LLVM 工具链中，覆盖率采集模块输出的覆盖率数据文件仅包含.profraw 文件，作用与 .gcda 相同，仅记录各代码行的执行次数，无法直接映射到源代码行中。因此，解析时需要结合构建包中的动态链接库及可重定位对象文件等确保覆盖率解析的准确性。这些文件提供了必要的调试符号和源代码映射信息。
        \item 在进行原始覆盖率数据处理时，LLVM 工具链下，.profraw 文件需先通过 llvm-profdata merge 合并，并结合构建包中的二进制文件生成可读的 .info 文件。而 GCC 工具链可以直接通过lcov工具，从 .gcda + .gcno 数据文件中提取覆盖率信息，生成 .info 文件。
    \end{itemize}
    
    基于上述差异，本系统在设计覆盖率解析流程时，明确区分两种工具链的处理方式，以确保能够使用通用的流程兼容两种不同的工具链。

    首先，用户将覆盖率数据包及对应的构建包上传至解析模块，其中构建包可选。解析模块对上传的压缩文件进行解压操作，并将解压后的文件保存至文件系统中以便后续处理。随后，解析模块查询数据库中保存的解析参数。上述LLVM工具链所需的二进制文件通常以压缩包形式存放在构建包中。因此此时的解析参数主要包括构建包中，二进制文件压缩包的路径，以及解压后的二进制文件路径。注意到某些领域较为特殊，需要对二进制文件进行额外处理，处理后二进制文件不再以压缩包形式存储，此时的解析参数仅有构建包中的二进制文件路径。

    在覆盖率解析的实际执行阶段，系统区分了两种常用编译工具链处理方式。对于LLVM工具链，解析模块首先从文件系统中读取原始覆盖率数据文件（.profraw），并调用 llvm-profdata 工具将这些原始数据合并为单一的中间文件（merged.raw），确保覆盖率统计的完整性与一致性。随后，解析模块根据解析参数读取相关的可执行二进制文件，并借助 llvm-cov 工具结合二进制文件从合并后的中间文件中导出覆盖率信息，生成标准化的 coverage.info 文件。对于 GCC 工具链，解析模块同样从文件系统中获取覆盖率数据文件（.gcda 与 .gcno），但直接通过 lcov 工具进行处理，并导出 coverage.info 文件。此过程保证了不同编译器生成的覆盖率数据能够统一解析，便于系统后续处理。

    生成 coverage.info 文件后，解析模块进一步调用 genhtml 工具，将 coverage.info 转化为可浏览的 HTML 格式覆盖率报告，同时保存到文件系统中，以便在前端系统进行可视化呈现。解析模块解析报告内容，提取关键覆盖率指标和统计信息，并将其结构化存储到数据库中，实现数据的持久化存储。整个流程体现了系统对覆盖率数据从上传、解压、工具调用、报告生成到数据库存储的端到端处理能力，并通过模块化设计与工具链解耦确保了系统的可扩展性与可维护性。

    注意到，上述覆盖率解析模块流程仅描述了全量覆盖率解析，而系统实际上实现了细粒度的覆盖率分析，能够按单板、领域、模块分层统计，支持增量覆盖率计算及多时间点对比和趋势分析。具体实现方面，主要包括模块覆盖率分析，以及增量覆盖率分析。

\subsubsection{模块覆盖率分析流程}
    模块覆盖率生成的过程视图如图~\ref{fig:4.8}所示，描述了从全量覆盖率中提取不同模块覆盖率数据的过程，可分为数据准备和模块覆盖解析两个阶段。

    \begin{figure}[ht]
        \centering
        \includegraphics[width=0.65\textwidth]{process/过程视图-覆盖率解析-模块覆盖率生成.pdf}
        \caption{过程视图-覆盖率解析-模块覆盖率生成流程时序图}
        \label{fig:4.8}
    \end{figure}

    在数据准备阶段，解析模块首先向数据库发起查询，获取用户配置的模块信息，其中包括模块名、模块与源文件的映射关系、模块的核心标记等。数据库返回查询结果后，解析模块对该信息进行解析，获取模块—文件映射，即每个模块包含的文件名列表，以及核心模块列表。同时，解析模块从文件系统中获取全量覆盖率文件（all.info）。

    接下来进入模块覆盖解析阶段。解析模块通过调用 lcov --extract 指令，根据模块包含的文件名列表，从 all.info 文件中提取与目标模块相关的覆盖率片段，生成对应的模块覆盖率数据文件（module.info）。进而，根据核心模块列表筛选出核心模块的 info 文件，利用 lcov add 指令将这些覆盖率数据合并，生成核心覆盖率文件（core.info），以便进行针对性分析。随后，解析模块基于各类 info 文件生成模块覆盖率报告。该报告以行粒度直观呈现覆盖率结果，并被存储至文件系统中，作为后续人工审阅与验证的依据。最后，解析模块对生成的模块覆盖率报告进行解析和结构化处理，提取其中的关键覆盖率指标，将其以模块为粒度存储至数据库，实现覆盖率数据的持久化管理。

    综上所述，该流程通过数据库、文件系统与解析模块之间的交互，有效实现了从全量覆盖率数据到模块覆盖率信息的提取、报告生成与持久化存储，确保覆盖率数据在不同抽象层次上的可追溯性和可管理性。这一设计不仅增强了覆盖率分析的针对性和精细化程度，同时为系统质量评估提供了坚实的数据支撑。

\subsubsection{增量覆盖率分析流程}
    增量覆盖率生成的过程视图如图~\ref{fig:4.9}所示，描述了从全量覆盖率中提取相对于基线的增量覆盖率数据的过程，可分为数据准备，计算增量代码和增量覆盖率解析两个阶段。

    \begin{figure}[ht]
        \centering
        \includegraphics[width=0.8\textwidth]{process/过程视图-覆盖率解析-增量覆盖率生成.pdf}
        \caption{过程视图-覆盖率解析-增量覆盖率生成流程时序图}
        \label{fig:4.9}
    \end{figure}

    在数据准备阶段，解析模块首先向数据库发起查询，获取当前节点信息和用户配置的基准节点信息。其中通过 git commit id 表示节点信息。同时，解析模块从文件系统中获取全量覆盖率文件（all.info）。

    在增量代码计算阶段，解析模块对 all.info 文件进行初步解析，通过字符串匹配，提取其中以“SF:”开头的行，过滤后即为覆盖率数据中包含的源文件路径列表。基于前缀匹配算法计算各源文件所在路径的最近公共祖先目录，从而确定待分析的代码范围。注意到，若此时所得祖先目录层级较高导致范围过大，则按照下一级目录进行拆分，分片进行后续操作再合并结果。确定了待分析目录后，解析模块调用 \verb|git diff <curr_commit> <base_commit> <path>| 命令，计算相应目录下，当前节点与基准节点之间的代码差异，并将差异结果保存为 dir.diff 文件，以便后续覆盖率比对使用。

    接下来进入增量覆盖率解析阶段，解析模块从文件系统中读取 dir.diff 与 all.info 文件。逐行解析 dir.diff 文件，识别出增量源代码行及其对应的函数信息。同时，解析模块结合 all.info 文件的内容，筛选出与上述增量代码行号和函数相匹配的覆盖率记录，从而生成增量覆盖率文件 diff.info。基于各 info 文件生成增量覆盖率报告。最后，解析模块对生成的增量覆盖率报告进行再次解析，提取其中的结构化覆盖率数据，并将其写入数据库，以支持后续的可视化展示与统计分析。
    
    综上所述，该流程通过数据库、文件系统与解析模块之间的交互，有效实现了从全量覆盖率数据到增量覆盖率信息的提取、报告生成与持久化存储，增强了覆盖率结果的可解释性与针对性，使覆盖率度量能够真正服务于快速迭代与高质量交付的工程目标。

\subsection{覆盖率可视化模块}
    覆盖率可视化模块的过程视图如图~\ref{fig:4.10}所示，描述了从用户请求发起，数据查询与处理，到最终结果呈现的完整的覆盖率可视化流程。

    \begin{figure}[ht]
        \centering
        \includegraphics[width=0.9\textwidth]{process/过程视图-覆盖率可视化流程.pdf}
        \caption{过程视图-覆盖率可视化流程时序图}
        \label{fig:4.10}
    \end{figure}

    当用户在界面中触发查看覆盖率数据的操作时，前端系统首先接收该请求，同时调用后端接口。后端系统在接收到请求后，依据调用参数访问数据库，完成对覆盖率数据的查询与过滤。数据库在完成数据检索后将结果返回至后端，由后端进一步进行业务逻辑层面的封装与处理，随后传递至前端。前端在得到数据后，还通过构造表格、柱状图、告警信息等多维度展示方式将其转化为用户可感知的可视化信息，提升了覆盖率信息的直观性和可解释性，最终呈现给用户。

    与此同时，系统还支持用户对覆盖率报告的进一步访问：用户可以点击柱状图中的柱子，查看对应的覆盖率报告。在该场景中，用户请求由前端解析并生成相应的 URL，向静态文件系统发起请求访问静态资源。静态文件系统返回存储的覆盖率报告文件，前端接收后直接渲染并展示给用户。此路径相较于动态查询更加高效，因为覆盖率报告作为静态资源已经完成预生成，无需额外的数据访问与计算。

    综上所述，该流程图揭示了覆盖率解析模块中从数据查询、逻辑处理到可视化生成的全链路过程，体现了系统在分层解耦、数据访问效率与用户体验之间的平衡。通过前后端协同与静态文件系统的引入，系统不仅能够保证覆盖率数据展示的实时性和准确性，还能够为用户提供结构化的覆盖率报告，以满足不同层级的分析与决策需求。

\section{物理视图}
    物理视图（Physical View）又称部署视图，面向系统部署，描述系统的交付、安装、部署的视图，主要解决系统安装部署的问题，描述系统的交付、安装、部署关系。 本节通过系统部署模型来描述部署视图。

    系统部署模型如图~\ref{fig:4.11}所示，主要包括四个物理节点：用户主机，目标硬件平台，应用服务器以及数据库服务器，同时引入 SmartCI 平台实现任务调度与自动化部署。

    \begin{itemize}
        \item \textbf{用户主机}：主要负责两部分功能：其一，冒烟测试组件，用于编写和提交测试用例，并触发测试流水线；其二，覆盖率看板，用网页形式呈现覆盖率可视化结果。用户只需在主机上完成测试用例的编写与提交，触发冒烟测试流水线后，其余步骤均由 SmartCI 平台自动化处理。
        \item \textbf{目标硬件平台}：分为控制面（LinuxOS）和数据面（异构OS），二者均包含测试用例执行环境和覆盖率采集组件。注意到，数据面缺乏独立文件系统且测试完成后自动断电，因此其覆盖率数据必须实时转发至控制面。在数据面，覆盖率采集组件通过插桩收集测试用例执行过程中的覆盖率信息，并利用共享内存通信传输到控制面。控制面的覆盖率采集组件汇总控制面自身和数据面的覆盖率信息，构造覆盖率数据文件，并通过 SmartCI 平台进行传输。
        \item \textbf{应用服务器}：运行 Flask 应用服务，包含覆盖率解析组件和静态文件系统。覆盖率解析组件负责将原始数据文件解析为结构化的覆盖率数据，同时生成 HTML 格式的覆盖率报告，分别存储于MySQL数据库和静态文件系统。Flask 服务通过统一接口访问数据库与静态文件系统，读取覆盖率数据和覆盖率报告生成前端页面供覆盖率看板访问，实现覆盖率可视化。
        \item \textbf{数据库服务器}：部署 MySQL 数据库，作为系统的数据持久化存储中心，用于存储覆盖率解析组件生成的数据。
        \item \textbf{SmartCI平台}：作为流水线任务调度中心，负责对用户提交的测试用例进行交叉编译，将可执行文件分发至目标硬件平台运行；测试完成后，自动收集覆盖率数据文件，并传输至应用服务器进行后续处理。该平台实现了任务的自动化与分布式执行，有效屏蔽了底层的任务调度与数据收集细节。
    \end{itemize}

    \begin{figure}[ht]
        \centering
        \includegraphics[width=0.9\textwidth]{physical/物理视图-部署模型.pdf}
        \caption{物理视图-部署模型图}
        \label{fig:4.11}
    \end{figure}

\section{本章小结}
    本章基于 4+1 视图方法对系统进行了概要设计。在需求分析阶段已通过用例图刻画了场景视图的基础上，本章进一步从逻辑视图、开发视图、过程视图与物理视图四个维度展开。逻辑视图侧重于系统功能模块及逻辑关系的分层建模；开发视图强调代码模型与实现组织的对应关系；过程视图通过运行时交互与时序机制揭示系统动态行为；物理视图则面向部署环境，明确系统在硬件与网络层面的映射关系。

    通过上述多维度的架构描述，本章在抽象层次上系统化呈现了覆盖率融合显示系统的整体设计思路。各视图既相互独立，又互为补充，从不同角度揭示了系统的结构、行为与实现机制，为后续的详细设计与系统实现奠定了坚实的架构基础。